\chapter{Differential Equations and Dynamical Systems}

The pendulum clock in the living room swings with an amplitude of sixty degrees, and the
length of its arm is one meter. The textbook gives the period
\[
T=2\pi\sqrt{L/g}=2.006409\ \text{seconds},
\]
and that formula is indeed correct for a very small swing. Measured at an amplitude of
sixty degrees, the same pendulum needs $2.153242$ seconds for one full swing: $7.32$
percent longer. A clock set with the small-angle formula falls behind by $5892$ seconds,
about ninety-eight minutes, every day.

The formula is not a miscalculation. It is the exact answer for a different equation.
The actual swing satisfies
\[
\frac{d^{2}\theta}{dt^{2}}=-\frac{g}{L}\sin\theta ,
\]
whereas the small-period formula comes from replacing $\sin\theta$ with $\theta$. At
$\theta=60^{\circ}$ that replacement turns $0.866025$ into $1.047198$, a difference of
seventeen percent at the farthest point, and the period error left behind is $7.32$
percent.

Chapter 12 taught us to differentiate, Chapter 14 series of functions, Chapter 15 vector
fields. Not one of them answers the question here: if the rate of change is known, how do we
find the function, is it unique, and what happens in the long run. This chapter fills those
four gaps: how to solve, conditions for existence and uniqueness, the phase plane, and the
stability test.

\begin{untukanak}
Imagine the water tank in a bathroom. The tap above runs at a steady rate, the drain below
runs out, and the higher the water, the faster it leaves. If we open the tap and then walk
away, tomorrow's water level cannot be computed from the size of the tap alone: what
determines it is the water level now. High water pushes out faster, the rise slows down,
and finally it stops when the water coming in equals the water going out.

Other things work the same way. Hot coffee cools faster while it is still very hot. Germs
multiply faster when there are many of them, then slow down when the food runs out. What we
know is the rule of change, not the final result.

Now the question that makes this work hard. Two tanks the same, taps the same, drains the
same, but one starts empty and one starts half full: do they end at the same level?
Sometimes yes, sometimes no. And there is a wicked drain, partly clogged and then suddenly
open: two different stories can come out of the same initial state.

In this chapter we shall learn why the rule of change sometimes suffices to determine the
whole future, and sometimes does not.
\end{untukanak}

\section{A rate that depends on itself}
\tingkat{basic}

A differential equation sets the rate of change of a function equal to something that
contains the function itself. To solve it means to find a function that satisfies that
sentence on an entire interval. The derivatives of Chapter 12 were computed at a single
point; the solution here is a function on an interval, and that difference is the one most
often missed.

\begin{definisi}[Ordinary differential equation]\index{differential equation}
The equation
\[
F\bigl(t,y,y',y'',\dots,y^{(n)}\bigr)=0
\]
is called an \emph{ordinary differential equation} (ODE) of order $n$ for a function $y$ of
the variable $t$ if $F$ contains the $n$-th derivative and contains no higher derivative.
The ODE is called \emph{linear} if $F$ has the form
\[
a_{n}(t)\,y^{(n)}+\cdots+a_{1}(t)\,y'+a_{0}(t)\,y=b(t),
\]
and is called \emph{autonomous} if $F$ does not contain $t$ explicitly.
\end{definisi}

 All the examples in
this chapter are of first or second order, because those two already contain every important
idea.

\begin{definisi}[Solution and initial value problem]\index{initial value problem}
A function $y$ is called a \emph{solution} of the ODE on an interval $I$ if $y$ has
derivatives up to order $n$ on $I$ and satisfies the equation at every point of $I$. If the
condition $y(t_{0})=y_{0}$ is added, the combination of the two is called an \emph{initial
value problem}, and $y_{0}$ is called the initial value.
\end{definisi}

One example to make sure the definition works. Take the population model
\[
y'=y(1-y),\qquad y(0)=\tfrac12 ,
\]
with $y$ denoting the fraction of the population relative to the carrying capacity. The
candidate solution is
\[
y(t)=\frac{1}{1+e^{-t}} .
\]
Its derivative is $y'(t)=e^{-t}/(1+e^{-t})^{2}$, and the right-hand side
$y(1-y)=e^{-t}/(1+e^{-t})^{2}$ is exactly the same. The initial value $y(0)=1/2$ matches as
well. So the function is a solution, and that check uses only the quotient rule of
Chapter 12.

What deserves attention: $y(t)=1$ and $y(t)=0$ are also solutions, both of them constant,
and both appear when we ask when the right-hand side is zero. So one equation can have many
solutions; the initial value selects one of them. The collection of all solutions is called
the \emph{general solution}, and it contains one constant for each order.

The interval on which a solution lives is part of the answer. Take
\[
y'=y^{2},\qquad y(0)=1 .
\]
Separation of variables gives $y(t)=1/(1-t)$, and that formula makes sense only on $t<1$.
At $t=0.9$ its value is $10$; at $t=0.99$ its value is $100$; at $t=0.999$ its value is
$1000$. There is no finite number at $t=1$, and no solution continues that function past it.
The solution \emph{blows up} in finite time, and a program that reports a value at $t=2$ does
not find a solution; it finds rounding.

\kamus{differential equation}{persamaan diferensial}{A sentence that ties the rate of change of a quantity to the size of that quantity itself.}

\kamus{general solution}{solusi umum}{The family of all functions satisfying the equation, usually containing one free constant per order.}

\begin{bacakeras}
--- If I let this ball go on the slope, it will move downhill. How fast?
--- It depends on where the ball is now. Near the top, slowly. In the middle of the slope,
quickly.
--- So its speed is not a single fixed number?
--- No. It is a rule: at each point, the speed is determined by that point.
--- If the rule already knows everything, is the path already determined?
--- Not yet. The rule gives a direction at every point; the path is born from the starting
point you choose.
--- Two balls, two starting points, one rule?
--- One rule, two paths.
--- If two paths meet at the same point, how do we know which one came from where?
--- We cannot. That is why this chapter demands one starting point for one path.

Repeat it once, in a normal voice: same rule, different starting point, different path.
\end{bacakeras}

\section{Two families that can be solved by hand}
\tingkat{basic}

Not every differential equation has a solution formula, but two large families do, and both
appear in almost every first-level application. The first family separates variables; the
second is the first-order linear equation.

\subsection{Separation of variables}

If the equation can be written $y'=f(t)g(y)$, we move all $y$ to one side and all $t$ to the
other side:
\[
\frac{dy}{g(y)}=f(t)\,dt,\qquad \int\frac{dy}{g(y)}=\int f(t)\,dt .
\]
That step is not a notational trick but the chain rule run backwards: if $H$ is an
antiderivative of $1/g$, then $\frac{d}{dt}H(y(t))=H'(y)y'=y'/g(y)$, so the original
equation becomes $\frac{d}{dt}H(y(t))=f(t)$ and Chapter 13 gives $H(y)=\int f\,dt+C$. The
result is usually implicit, and that is a legitimate answer.

The logistic model above is used throughout the chapter. With $g(y)=y(1-y)$ and partial
fractions,
\[
\int\frac{dy}{y(1-y)}=\int\Bigl(\frac{1}{y}+\frac{1}{1-y}\Bigr)dy
=\ln\Bigl|\frac{y}{1-y}\Bigr| ,
\]
so $\ln|y/(1-y)|=t+C$ and
\[
y(t)=\frac{1}{1+Ae^{-t}},\qquad A=\frac{1-y_{0}}{y_{0}} .
\]
For $y_{0}=1/2$ we obtain $A=1$, which is the formula we checked a moment ago. Note that $A$
is determined by the initial value, and that $A=0$ gives the constant solution $y=1$ which
the division above does not cover: when separating variables we divide by $g(y_{0})$, and
that division is valid only if $g(y_{0})\ne0$.

\subsection{Linear equations and the integrating factor}

The standard form of a first-order linear equation is
\[
y'+p(t)\,y=q(t).
\]
The left-hand side is almost the derivative of a product, only one factor short. The missing
factor is exactly what the following theorem looks for.

\begin{teorema}[Integrating factor]\index{integrating factor}
Suppose $p$ and $q$ are continuous on an open interval $I$ containing $t_{0}$, and suppose
\[
\mu(t)=\exp\Bigl(\int_{t_{0}}^{t}p(s)\,ds\Bigr).
\]
Then the initial value problem $y'+p(t)y=q(t)$, $y(t_{0})=y_{0}$ has exactly one solution on
$I$, namely
\[
y(t)=\frac{1}{\mu(t)}\Bigl(y_{0}+\int_{t_{0}}^{t}\mu(s)\,q(s)\,ds\Bigr).
\]
\end{teorema}

\begin{proof}
Since $p$ is continuous, $\mu$ is defined, positive, and differentiable on $I$, with
$\mu'=p\mu$ by the fundamental theorem of calculus (Chapter 13). Multiply the equation by
$\mu$:
\[
\mu y'+\mu p\,y=\mu q
\quad\Longleftrightarrow\quad
(\mu y)'=\mu q ,
\]
because the product rule gives $(\mu y)'=\mu'y+\mu y'=\mu(py+y')$. Integrate both sides from
$t_{0}$ to $t$ and use the fundamental theorem of calculus:
\[
\mu(t)y(t)-\mu(t_{0})y(t_{0})=\int_{t_{0}}^{t}\mu(s)q(s)\,ds .
\]
Since $\mu(t_{0})=1$, division by $\mu(t)>0$ gives the formula requested. Uniqueness: if
$y_{1},y_{2}$ are two solutions, their difference $w=y_{1}-y_{2}$ satisfies $w'+pw=0$ with
$w(t_{0})=0$; multiplying by $\mu$ gives $(\mu w)'=0$, so $\mu w$ is constant, and since its
value is zero at $t_{0}$ it is zero everywhere.
\end{proof}

That proof guesses nothing: $\mu$ is found by demanding $\mu'=p\mu$, an equation easier than
the original. This pattern repeats: one problem is solved by turning it into a simpler
problem.

\begin{contoh}[One linear equation, two checks]
Solve $y'+2y=4t$ with $y(0)=1$. Here $p=2$ and $q=4t$, so $\mu=e^{2t}$ and
\[
(e^{2t}y)'=4t\,e^{2t}.
\]
Integration by parts gives
\[
\int 4t\,e^{2t}\,dt=\Bigl(2t-1\Bigr)e^{2t}+C ,
\]
because the derivative of $(2t-1)e^{2t}$ is $2e^{2t}+2(2t-1)e^{2t}=(4t)e^{2t}$. Thus
\[
y(t)=2t-1+Ce^{-2t},
\]
and the condition $y(0)=1$ gives $C=2$, that is $y(t)=2t-1+2e^{-2t}$.

Check: substitution gives $y'+2y=4t$, and $y(1)=1+2e^{-2}=1.2706705665$; an adaptive
integrator with tolerance $10^{-12}$ gives $1.270670566473$.
\end{contoh}

\subsection{The RC circuit: the simplest real example}

A voltage source $V_{s}$ is connected in series with a resistor $R$ and a capacitor $C$.
Kirchhoff's law of voltages gives $V_{s}=Ri+V_{C}$, and the law of the capacitor gives
$i=C\,dV_{C}/dt$. The two together give
\[
RC\,\frac{dV_{C}}{dt}+V_{C}=V_{s},
\qquad\text{or}\qquad
\frac{dV_{C}}{dt}+\frac{1}{RC}V_{C}=\frac{V_{s}}{RC}.
\]
This is a linear equation with $p=1/(RC)$ constant. By the theorem above, with
$V_{C}(0)=0$,
\[
V_{C}(t)=V_{s}\bigl(1-e^{-t/\tau}\bigr),\qquad \tau=RC .
\]
The number $\tau$ is called the time constant. For $R=1\,\mathrm{k}\Omega$ and
$C=1\,\mu\mathrm{F}$, $\tau=10^{-3}$ seconds. With $V_{s}=5$ volts, the voltage at
$t=\tau$ is $3.160603$ volts and at $t=5\tau$ it is $4.966310$ volts, that is $99.33$
percent of $5$ volts. To reach $90$ percent requires $t=\tau\ln10=2.302585$ milliseconds.

\begin{resep}
\begin{enumerate}
\item Write the equation in standard form $y'+p(t)y=q(t)$: divide by the coefficient of
  $y'$ so that the coefficient is exactly one.
\item Compute $\mu(t)=\exp\bigl(\int p(t)\,dt\bigr)$, with no constant of integration.
\item Multiply the whole equation by $\mu$, then write the left-hand side as $(\mu y)'$. If
  the left-hand side refuses to become a single derivative, step 1 is wrong.
\item Integrate both sides with respect to $t$, then divide by $\mu$.
\item Use the initial value to determine the constant, and write down the interval of
  validity.
\end{enumerate}
Warning: step 1 must not be skipped, and step 3 must not be done before step 2. Computing
$\mu$ from coefficients that have not yet been normalized gives a wrong integrating factor,
and the error only becomes visible when the solution is tested.
\end{resep}

\kamus{integrating factor}{faktor integrasi}{A multiplier that turns the left-hand side of a linear equation into a single derivative of a product.}

\section{Existence, uniqueness, and Picard iteration}
\tingkat{intermediate}

The example $y'=y^{2}$ shows that a solution can stop existing. Given $f$ and one initial
point, does the solution exist and is it only one? The answer depends on one property of $f$
with respect to $y$.

\begin{definisi}[Lipschitz condition]\index{Lipschitz condition}
A function $f(t,y)$ satisfies the Lipschitz condition with respect to $y$ on a region
$D\subseteq\mathbb{R}^{2}$ if there is a constant $L$ such that
\[
|f(t,y_{1})-f(t,y_{2})|\le L\,|y_{1}-y_{2}|
\]
for every $(t,y_{1}),(t,y_{2})\in D$. The smallest possible constant is called the Lipschitz
constant.
\end{definisi}

Its everyday sentence: $f$ must not change too abruptly when $y$ shifts a little. If $f$ has
a partial derivative with respect to $y$ that is continuous and bounded by $L$ on $D$, that
condition holds automatically, because the mean value theorem (Chapter 12) gives
$f(t,y_{1})-f(t,y_{2})=\partial_{y}f(t,\xi)(y_{1}-y_{2})$.

\begin{teorema}[Picard--Lindel\"of]\index{Picard-Lindelof theorem}
Suppose $f$ is continuous on the closed rectangle
$D=[t_{0}-a,t_{0}+a]\times[y_{0}-b,y_{0}+b]$ and satisfies the Lipschitz condition with
respect to $y$ on $D$ with constant $L$. Suppose $M=\max_{D}|f|$ and
$h=\min(a,\,b/M)$. Then the initial value problem $y'=f(t,y)$, $y(t_{0})=y_{0}$ has exactly
one solution on the interval $|t-t_{0}|\le h$.
\end{teorema}

The full proof uses the Banach fixed point theorem on the space of continuous functions with
the maximum norm, which is complete because uniform continuity is discussed in Chapter 11
and Chapter 14. I give the sketch, because the sketch contains the whole idea and can be
turned into an algorithm.

Write the equation as an integral equation, which is equivalent because both sides are
continuous:
\[
y(t)=y_{0}+\int_{t_{0}}^{t}f\bigl(s,y(s)\bigr)\,ds .
\]
Define the map $T$ by that formula, starting from the constant $y_{0}(t)=y_{0}$, then
iterate:
\[
y_{n+1}(t)=y_{0}+\int_{t_{0}}^{t}f\bigl(s,y_{n}(s)\bigr)\,ds .
\]
Every $y_{n}$ lies inside $D$ as long as $|t-t_{0}|\le h$, because
$|y_{n+1}-y_{0}|\le M|t-t_{0}|\le Mh\le b$. The difference of two successive iterations
satisfies
\[
|y_{n+1}(t)-y_{n}(t)|\le L\Bigl|\int_{t_{0}}^{t}|y_{n}(s)-y_{n-1}(s)|\,ds\Bigr|
\le \frac{(L|t-t_{0}|)^{n}}{n!}\,\|y_{1}-y_{0}\| ,
\]
and induction on $n$ gives that bound. Since $\sum_{n}(Lh)^{n}/n!$ converges
(ratio test, Chapter 5), the series $y_{0}+\sum(y_{n+1}-y_{n})$ converges uniformly, and
its limit is the solution. Uniqueness comes from the same inequality: two different solutions
would satisfy $|y_{1}-y_{2}|\le L|t-t_{0}|\,\|y_{1}-y_{2}\|$, which for $Lh<1$ is possible
only if the difference is zero; the case $Lh\ge1$ is handled by shrinking $h$, and that is
why the theorem promises only a short interval.

Run that iteration on $y'=y$, $y(0)=1$. Starting from $y_{0}=1$:
\[
y_{1}=1+t,\quad y_{2}=1+t+\frac{t^{2}}{2},\quad
y_{3}=1+t+\frac{t^{2}}{2}+\frac{t^{3}}{6},
\]
so $y_{n}$ is the partial sum of the exponential series from Chapter 5. At $t=1$, the eighth
iteration gives $2.718278769841$, whereas $e=2.718281828459$; the error is
$3.06\times10^{-6}$.

The Lipschitz condition must not be thrown away. Take
\[
y'=\sqrt{|y|},\qquad y(0)=0 .
\]
The right-hand side is continuous, so a solution exists; but it does not satisfy the
Lipschitz condition near $y=0$, because its slope is unbounded there. And indeed there are
two solutions: the zero function, $y(t)\equiv0$, and
\[
y(t)=\frac{t^{2}}{4}\quad (t\ge0).
\]
The check is one line: $y'=t/2$ and $\sqrt{|y|}=\sqrt{t^{2}/4}=t/2$ for
$t\ge0$, whereas $y(0)=0$. Two different histories from one and the same initial state,
exactly like the drain that opens suddenly in the opening box of the chapter.

Lipschitz is a \emph{sufficient} condition, not a necessary one. Existence without
uniqueness is guaranteed by Peano's theorem, which uses compactness (Chapter 11) and
Arzel\`a--Ascoli; see Coddington and Levinson, chapters 1--2.

\kamus{uniqueness}{ketunggalan}{The guarantee that one initial state gives birth to only one history, with no branch.}
\section{Systems, the phase plane, and oscillators}
\tingkat{intermediate}

One second-order equation is not a new animal; it is two first-order equations running
together. By naming $v=\theta'$, the pendulum equation becomes
\[
\theta'=v,\qquad v'=-\frac{g}{L}\sin\theta .
\]
This pair is autonomous: its right-hand side does not contain $t$, so the initial state
determines the whole history (provided the Lipschitz condition holds) and we may forget time
for a moment in order to ask what happens on the plane $(\theta,v)$.

\begin{definisi}[Autonomous systems and the phase plane]\index{phase plane}
The system
\[
\mathbf{x}'=\mathbf{f}(\mathbf{x}),\qquad \mathbf{x}\in\mathbb{R}^{n},
\]
is called \emph{autonomous}. For $n=2$, the plane on which the point $\mathbf{x}$ lives is
called the \emph{phase plane}; the curve $\{\mathbf{x}(t)\}$ is called a
\emph{trajectory}, and the picture of all the arrows $\mathbf{f}(\mathbf{x})$ is called the
\emph{direction field}. A point $\mathbf{x}^{*}$ with $\mathbf{f}(\mathbf{x}^{*})=0$ is
called an \emph{equilibrium point} or a fixed point.
\end{definisi}

The horizontal axis of the phase plane represents position, the vertical axis represents
velocity, and one point means ``here, with this speed''. Time is not drawn; it is only the
parameter that walks along the curve. That is its advantage: the whole history, for all
initial points at once, appears as a single picture.

The equilibrium points of the pendulum are easy to find. $\theta'=0$ demands $v=0$, and
$v'=0$ demands $\sin\theta=0$, so $\theta=k\pi$. At $\theta=0$ the pendulum hangs straight
down; at $\theta=\pi$ it stands straight up. The first is stable in the everyday sense, the
second is not: the slightest touch brings it down.

\begin{untukanak}
Go to the park, find a swing, and sit down. Ask a friend to push you gently, only a little.
Count ten swings by counting inside your head: one, two, three, up to ten. Now ask for a push
that is much higher, until the swing rises almost as high as the upper bar, then count ten
swings again.

What you will see: the high swing takes longer for one full round, even though you feel as
though the high swing is faster. That feeling cheats because the distance traveled is much
longer, while the time for one round is almost the same.

Now look at the top of the swing. When you are at the highest point, there is a short moment
when you almost stop and hang. At the lowest point, you move fastest. A swing never stops in
the middle, and is never fastest at the ends. Those two things always change places.

In this chapter we shall learn why a high swing is slower, and why a pendulum standing
straight up is almost impossible to keep there.
\end{untukanak}

Now the exact period formula. Multiply $v'=-(g/L)\sin\theta$ by $v=\theta'$
and integrate:
\[
\frac{d}{dt}\Bigl(\frac{v^{2}}{2}+\frac{g}{L}(1-\cos\theta)\Bigr)=0 .
\]
That bracket is the energy per unit mass, and it is constant along the trajectory.
With $E$ denoting its value,
\[
\frac{v^{2}}{2}+\frac{g}{L}\bigl(1-\cos\theta\bigr)=E ,
\]
so $v=\pm\sqrt{2E-2(g/L)(1-\cos\theta)}$. The trajectories in the phase plane are those
energy curves; in Figure~\ref{fig:bidang-fase-pendulum} they are drawn for the unit
$g/L=1$.

\begin{figure}[htbp]
\centering
\begin{tikzpicture}[x=1.05cm,y=1.0cm,font=\footnotesize]
  \draw[->,layergrey] (-4.15,0) -- (4.15,0) node[right] {$\theta$};
  \draw[->,layergrey] (0,-2.95) -- (0,2.95) node[above] {$v$};
  \foreach \tick/\lab in {-3.1416/{$-\pi$},-1.5708/{$-\pi/2$},
                          1.5708/{$\pi/2$},3.1416/{$\pi$}}{
    \draw[layergrey!55] (\tick,0.07) -- (\tick,-0.07);
    \node[layergrey,font=\scriptsize,below] at (\tick,-0.07) {\lab};}
  \node[layergrey,font=\scriptsize,below left] at (-0.06,0.02) {$0$};
  \foreach \th in {-3.5,-2.6,-1.7,-0.85,0.85,1.7,2.6,3.5}{
    \foreach \w in {-2.4,-1.6,-0.8,0,0.8,1.6,2.4}{
      \pgfmathsetmacro{\nn}{sqrt(\w*\w+sin(\th r)*sin(\th r))}
      \pgfmathsetmacro{\dx}{0.32*\w/\nn}
      \pgfmathsetmacro{\dy}{-0.32*sin(\th r)/\nn}
      \draw[-{Latex[length=0.9mm]},layerteal!60]
        (\th,\w) -- ({\th+\dx},{\w+\dy});}}
  \foreach \En/\tm in {0.4/0.9273,1.0/1.5708,1.6/2.2143}{
    \draw[layerindigo,thick]
      plot[domain=-\tm:\tm,samples=121,variable=\t]
        ({\t},{sqrt(max(0,2*(\En-1+cos(\t r))))});
    \draw[layerindigo,thick]
      plot[domain=-\tm:\tm,samples=121,variable=\t]
        ({\t},{-sqrt(max(0,2*(\En-1+cos(\t r))))});}
  \draw[accent,very thick]
    plot[domain=-3.1416:3.1416,samples=161,variable=\t]
      ({\t},{2*cos(\t r/2)});
  \draw[accent,very thick]
    plot[domain=-3.1416:3.1416,samples=161,variable=\t]
      ({\t},{-2*cos(\t r/2)});
  \draw[layerplum,thick,densely dashed]
    plot[domain=-3.6:3.6,samples=161,variable=\t]
      ({\t},{sqrt(2*(1.6+cos(\t r)))});
  \fill[ink] (0,0) circle (1.6pt);
  \node[ink,anchor=north east,font=\scriptsize] at (-0.12,-0.52) {center};
  \fill[accent] (3.1416,0) circle (1.6pt);
  \fill[accent] (-3.1416,0) circle (1.6pt);
  \node[accent,anchor=north west,font=\scriptsize,align=left]
    at (2.35,-0.18) {saddle};
  \node[accent,anchor=west,font=\scriptsize,align=left]
    at (2.30,1.30) {separatrix $E=2$};
  \node[layerplum,anchor=west,font=\scriptsize,align=left]
    at (0.55,2.72) {open trajectory $E=2.6$};
\end{tikzpicture}
\caption{The phase plane of the pendulum in units $g/L=1$, with the horizontal axis the angle
$\theta$ and the vertical axis the velocity $v$. The small arrows show the direction field
$(\theta',v')=(v,-\sin\theta)$. The dark blue curves are three closed trajectories with
energy $E=0.4$, $1.0$, and $1.6$; the larger $E$, the wider the turn. The red curve is
the separatrix $E=2$, the trajectory that reaches exactly $\theta=\pm\pi$ with zero velocity;
it separates closed swings from open rotations. The dashed curve is the open trajectory
$E=2.6$, where the pendulum rotates fully without ever stopping. The point at $\theta=0$ is
a center, the points at $\theta=\pm\pi$ are saddles.}
\label{fig:bidang-fase-pendulum}
\end{figure}

That figure contains the entire qualitative behavior of the pendulum without a single
equation being solved; all that is needed is the sign of $\sin\theta$ and the symmetry of the
energy.

For small amplitude, $\sin\theta\approx\theta$ and the system becomes the harmonic
oscillator
\[
\theta'=v,\qquad v'=-\omega_{0}^{2}\theta,\qquad \omega_{0}=\sqrt{g/L},
\]
with solution $\theta(t)=\theta_{0}\cos(\omega_{0}t)$ and period
$2\pi/\omega_{0}=2.006409$ seconds for $L=1$ meter. Its trajectories are the ellipses
$v^{2}+\omega_{0}^{2}\theta^{2}=\text{constant}$, and those are the three small rings near
the center in the figure above.

The exact period comes from energy conservation, not from linearization.
With $E=(g/L)(1-\cos\theta_{0})$ for a pendulum released from rest at
$\theta_{0}$, the time to fall from $\theta_{0}$ to zero is
\[
\frac{T}{4}=\sqrt{\frac{L}{2g}}\int_{0}^{\theta_{0}}
\frac{d\theta}{\sqrt{\cos\theta-\cos\theta_{0}}}
=\sqrt{\frac{L}{g}}\,K(k),\qquad k=\sin\frac{\theta_{0}}{2},
\]
with $K$ the complete elliptic integral of the first kind. For $\theta_{0}=60^{\circ}$
we obtain $k=0.5$ and $K$ at the parameter $m=k^{2}=0.25$ has the value
$1.6857503548$, so $T=2.153242$ seconds. Quadrature on that integral gives the same
number to within $9\times10^{-13}$; two different roads, one result. For $\theta_{0}=5^{\circ}$
the ratio is $1.00047617$, that is an error of less than one part in two thousand: that is
where the small formula deserves to be used, and where its limit of validity lies.

\begin{tangan}
Take a string one meter long, a metal nut as the weight, and a stopwatch on the phone. Tie
the nut to the end of the string, stick the other end to the door frame with tape, and
measure the length of the string from the suspension point to the center of the nut: try to
make it exactly one meter. Pull the nut aside only about five centimeters, release it, and
record the time for ten full swings. Divide by ten; the result is about two seconds.

Now pull much higher, roughly half a meter above the lowest point, and repeat. Compare the
two numbers. What you will see: the second number is larger, but not twice as large. The
difference is about seven percent, and that agrees with the computation above if the angle is
about sixty degrees.

Repeat with a string a quarter of a meter long: the period becomes half, because the period
depends on the square root of the length.
\end{tangan}

\section{Stability: eigenvalues, linearization, and Lyapunov}
\tingkat{advanced}

The question most often asked about a system is not ``what is its value at $t=7$'' but ``does it
come back after being disturbed''. To answer it we do not need the solution; the derivatives
around the equilibrium point are enough.

\begin{definisi}[Stability]\index{stability}
Let $\mathbf{x}^{*}$ be an equilibrium point of $\mathbf{x}'=\mathbf{f}(\mathbf{x})$.
That point is called \emph{stable} if for every $\varepsilon>0$ there is
$\delta>0$ such that $\|\mathbf{x}(0)-\mathbf{x}^{*}\|<\delta$ guarantees
$\|\mathbf{x}(t)-\mathbf{x}^{*}\|<\varepsilon$ for every $t\ge0$. The point is
called \emph{asymptotically stable} if it is stable and there is $\delta>0$ such that
$\mathbf{x}(t)\to\mathbf{x}^{*}$ as $t\to\infty$. If it is not stable, it is
called \emph{unstable}.
\end{definisi}

\begin{teorema}[Linearization]\index{linearization}
Suppose $\mathbf{f}$ has continuous partial derivatives near the equilibrium point
$\mathbf{x}^{*}$ and suppose $A=J_{\mathbf{f}}(\mathbf{x}^{*})$ is the Jacobian matrix from
Chapter 12. If no eigenvalue of $A$ has real part zero (a \emph{hyperbolic} equilibrium
point), then the qualitative behavior of the nonlinear system near $\mathbf{x}^{*}$ is the
same as that of the linear system $\mathbf{u}'=A\mathbf{u}$: the type of stability is the
same, and nearby trajectories can be mapped one-to-one and continuously onto linear
trajectories.
\end{teorema}

This is the Hartman--Grobman theorem; its proof is long and uses the fixed point theorem on a
space of functions, so I cite it: Hartman, chapter 9, Perko, section 2.8. The practical
consequence: local stability is read off from the eigenvalues.

For a two-dimensional system, every possibility is read from the trace
$\tau=\operatorname{tr}A$ and the determinant $\Delta=\det A$; the eigenvalues are the roots
of $\lambda^{2}-\tau\lambda+\Delta=0$. If $\Delta<0$, the roots are real and of
opposite sign: a \emph{saddle}. If $\Delta>0$ and $\tau^{2}\ge4\Delta$,
both are real and of the same sign: a \emph{node}, stable if $\tau<0$. If
$\tau^{2}<4\Delta$, the roots are conjugates with real part $\tau/2$:
a \emph{spiral}, stable if $\tau<0$, and a center if $\tau=0$.

Example: the matrix $A=\begin{pmatrix}0&1\\-2&-3\end{pmatrix}$ has $\tau=-3$
and $\Delta=2$, so its roots are $\lambda^{2}+3\lambda+2=0$, namely $-1$ and
$-2$: a stable node. The matrix $B=\begin{pmatrix}0&1\\-2&0\end{pmatrix}$ has
$\tau=0$ and $\Delta=2>0$, its roots are $\pm i\sqrt{2}$: a center.

Figure~\ref{fig:potret-linear} shows the two portraits that differ most in character: the
saddle, where almost every initial point moves away, and the stable spiral, where everything
spirals inward.

\begin{figure}[htbp]
\centering
\begin{tikzpicture}[font=\footnotesize]
  \begin{scope}[x=1.35cm,y=1.35cm]
    \begin{scope}
      \clip (-1.62,-1.52) rectangle (1.62,1.52);
      \foreach \c in {0.05,0.30,0.70,1.20}{
        \draw[layerteal,thick] plot[domain=-1.55:1.55,samples=90,variable=\s]
          ({sqrt(\c+\s*\s)},{\s});
        \draw[layerteal,thick] plot[domain=-1.55:1.55,samples=90,variable=\s]
          ({-sqrt(\c+\s*\s)},{\s});}
      \draw[layerindigo,densely dashed] (-1.5,-1.5) -- (1.5,1.5);
      \draw[layerindigo,densely dashed] (-1.5,1.5) -- (1.5,-1.5);
    \end{scope}
    \draw[->,layergrey] (-1.75,0) -- (1.75,0) node[right] {$x$};
    \draw[->,layergrey] (0,-1.55) -- (0,1.55) node[above] {$y$};
    \fill[ink] (0,0) circle (1.5pt);
    \node[layerindigo,font=\scriptsize,anchor=south east] at (-1.52,1.30)
      {eigen directions};
    \node[layergrey,anchor=north] at (0,-1.62)
      {(a) saddle: $\lambda=1,-1$};
  \end{scope}
  \begin{scope}[xshift=6.6cm,x=1.35cm,y=1.35cm]
    \draw[->,layergrey] (-1.75,0) -- (1.75,0) node[right] {$x$};
    \draw[->,layergrey] (0,-1.55) -- (0,1.55) node[above] {$y$};
    \foreach \r in {0.35,0.7,1.1,1.5}{
      \draw[accent,thick]
        plot[domain=0:12.6,samples=200,variable=\t]
          ({\r*exp(-0.3*\t)*cos(\t r)},{\r*exp(-0.3*\t)*(-sin(\t r))});}
    \fill[ink] (0,0) circle (1.5pt);
    \node[layergrey,anchor=north] at (0,-1.62)
      {(b) stable spiral: $\lambda=-0.3\pm i$};
  \end{scope}
\end{tikzpicture}
\caption{Two linear phase portraits of different types. (a) The system
$x'=y$, $y'=x$: trajectories satisfy $x^{2}-y^{2}=\text{constant}$, hyperbolas
with asymptotes $y=\pm x$, and the dashed lines are the eigen directions, one in
and one out, so the origin is a saddle and is unstable. (b) The system $x'=-0.3x+y$, $y'=-x-0.3y$:
the eigenvalues are $-0.3\pm i$, so trajectories spiral while shrinking by the
factor $e^{-0.3t}$ and the origin is a stable spiral. The four trajectories in (b)
differ only in their initial radius, because the system is linear.}
\label{fig:potret-linear}
\end{figure}

One warning that is expensive to learn. If $\tau=0$ and $\Delta>0$, linearization gives a
center, but the nonlinear system at that point is not necessarily a center. Take
\[
x'=-y+x(x^{2}+y^{2}),\qquad y'=x+y(x^{2}+y^{2}).
\]
Its Jacobian at the origin is a rotation matrix, its eigenvalues are $\pm i$:
linearization says center. But in polar coordinates $r^{2}=x^{2}+y^{2}$,
the cross multiplication gives
\[
\frac{1}{2}\frac{d}{dt}r^{2}=xx'+yy'=r^{2}(x^{2}+y^{2})=r^{4},
\]
so $r'=r^{3}>0$ for $r>0$: trajectories spiral outward, and the origin is not
stable. Linearization is not wrong; it is only insufficient when there is an eigenvalue on
the imaginary axis. That is where the theorem above demands a hyperbolic equilibrium point.

For the non-hyperbolic case, a second tool is used: a decreasing energy function, which we
compose ourselves.

\begin{definisi}[Lyapunov function]\index{Lyapunov function}
A function $V$ that is continuous and has continuous derivatives on a neighborhood $U$ of
$\mathbf{x}^{*}$ is called a \emph{Lyapunov function} if $V(\mathbf{x}^{*})=0$ and
$V(\mathbf{x})>0$ for $\mathbf{x}\in U\setminus\{\mathbf{x}^{*}\}$. Its derivative along a
trajectory is written
\[
\dot V(\mathbf{x})=\nabla V(\mathbf{x})\cdot\mathbf{f}(\mathbf{x}).
\]
\end{definisi}

That artificial hill is $V$, and $\dot V$ is the steepness of its descent.

\begin{teorema}[Lyapunov test]
Suppose $V$ is a Lyapunov function near the equilibrium point $\mathbf{x}^{*}$.
\begin{enumerate}
\item If $\dot V\le0$ on $U$, then $\mathbf{x}^{*}$ is stable.
\item If $\dot V<0$ on $U\setminus\{\mathbf{x}^{*}\}$, then
  $\mathbf{x}^{*}$ is asymptotically stable.
\end{enumerate}
\end{teorema}

\begin{proof}
For the first statement, take $\varepsilon>0$ small enough that the ball
$\|\mathbf{x}-\mathbf{x}^{*}\|\le\varepsilon$ stays inside $U$. On that ball
$V$ is positive and the ball is compact, so
$m=\min_{\|\mathbf{x}-\mathbf{x}^{*}\|=\varepsilon}V>0$. Choose $\delta>0$
such that $V(\mathbf{x})<m$ for $\|\mathbf{x}-\mathbf{x}^{*}\|<\delta$;
this is possible because $V$ is continuous and $V(\mathbf{x}^{*})=0$. A trajectory starting
in the $\delta$ ball satisfies $V(\mathbf{x}(t))\le V(\mathbf{x}(0))<m$, since $\dot V\le0$,
so it cannot cross the $\varepsilon$ ball; hence it is stable.

For the second statement, $V(\mathbf{x}(t))$ decreases and is bounded below by
zero, so it has a limit $L\ge0$. Suppose $L>0$; the trajectory then stays in the compact set
$\{V\ge L\}$, and there $-\dot V$ has a positive minimum, so $V$ would have to fall through
$L$ in finite time. Contradiction. Hence $L=0$, and because $V$ vanishes only at
$\mathbf{x}^{*}$, the trajectory heads there.
\end{proof}

An example is the damped pendulum, $\theta'=v$, $v'=-(g/L)\sin\theta-bv$ with
$b>0$, and $V=v^{2}/2+(g/L)(1-\cos\theta)$, which is the mechanical energy, positive
except at the equilibrium point. Its derivative is
\[
\dot V=v\Bigl(-\frac{g}{L}\sin\theta-bv\Bigr)+\frac{g}{L}\sin\theta\,v=-bv^{2}\le0 .
\]
So the lower equilibrium point is stable and the energy never increases. This test
does not give asymptotic stability, because $\dot V$ is zero along the whole line $v=0$;
the LaSalle version adds that the only invariant set inside $\{\dot V=0\}$ is the equilibrium
point, so all trajectories head there.

Linearization gives the same answer. At $\theta=0$ the Jacobian is
$\begin{pmatrix}0&1\\-g/L&-b\end{pmatrix}$, with $\tau=-b$ and
$\Delta=g/L>0$; for $b=0.5$ and $g/L=9.80665$ the roots are $-0.25\pm3.121562\,i$,
that is a stable spiral. At $\theta=\pi$ a $+g/L$ appears in the lower right corner,
so $\Delta=-g/L<0$: a saddle with eigenvalues $\pm3.131557$. A pendulum standing
straight up is indeed a saddle, and not as a figure of speech: there is exactly one direction
that keeps it up there, and every disturbance carries it away.

\kamus{asymptotically stable}{stabil asimtotik}{Stays close after a disturbance, then actually arrives back at the equilibrium point.}

\kamus{Lyapunov function}{fungsi Lyapunov}{An artificial hill used as a ruler: as long as the system descends it, the system cannot run far away.}

\section{Euler, Runge--Kutta, and stiffness}
\tingkat{advanced}

Almost no equation can be solved by a formula, so numerical methods are not a supplement but
the main road. The idea goes back to Chapter 12: if we know the value and the derivative at
one point, we can estimate the value a little to its right.

Euler's method does the simplest thing. From $y'=f(t,y)$ and the initial value,
\[
y_{n+1}=y_{n}+h\,f(t_{n},y_{n}),\qquad t_{n+1}=t_{n}+h .
\]
It is exactly the first two terms of the Taylor series around $t_{n}$, so its one-step error
is the term $h^{2}y''(\xi)/2$, of order $h^{2}$. That local error accumulates into a global
error of order $h$: doubling the accuracy demands twice as many steps, not four times.

Numbers for $y'=y$, $y(0)=1$, at $t=1$ with $h=0.1$: Euler gives
$(1.1)^{10}=2.5937424601$, whereas $e=2.7182818285$. The error is $0.1245$, that is
$4.58$ percent, and the error of the first single step is $5.171\times10^{-3}$.

The fourth-order Runge--Kutta method repairs this by taking four samples of the slope at
every step:
\[
y_{n+1}=y_{n}+\frac{h}{6}\bigl(k_{1}+2k_{2}+2k_{3}+k_{4}\bigr),
\]
with $k_{1}=f(t_{n},y_{n})$, $k_{2}=f(t_{n}+h/2,\,y_{n}+hk_{1}/2)$,
$k_{3}=f(t_{n}+h/2,\,y_{n}+hk_{2}/2)$, and $k_{4}=f(t_{n}+h,\,y_{n}+hk_{3})$.
Its local error is of order $h^{5}$ and its global error of order $h^{4}$: half the step
means the error falls by about sixteen times.

The following table comes from the program in the lab section, on the logistic model with
$y(0)=1/2$ and $t=2$; the exact value is $0.8807970780$.

\begin{center}
{\small
\begin{tabular}{lrrr}
\hline
$h$ & Euler error & RK4 error & Euler ratio \\
\hline
$0.400$ & $1.958\times10^{-2}$ & $1.245\times10^{-5}$ & --- \\
$0.200$ & $9.423\times10^{-3}$ & $7.237\times10^{-7}$ & $2.08$ \\
$0.100$ & $4.630\times10^{-3}$ & $4.359\times10^{-8}$ & $2.04$ \\
$0.050$ & $2.296\times10^{-3}$ & $2.674\times10^{-9}$ & $2.02$ \\
$0.025$ & $1.143\times10^{-3}$ & $1.656\times10^{-10}$ & $2.01$ \\
\hline
\end{tabular}}
\end{center}

The last column is the proof of first order: every time $h$ is halved, the Euler error falls
by almost exactly two; for RK4 the ratio is about $16$, in keeping with fourth order.
Figure~\ref{fig:galat-numerik} displays the same numbers on a logarithmic scale on both
axes.

\begin{figure}[htbp]
\centering
\begin{tikzpicture}
\begin{axis}[width=0.80\linewidth,height=5.6cm,
  xlabel={step $h$},ylabel={$|y(2)-y_{h}(2)|$},
  xmode=log,ymode=log,xmin=0.02,xmax=0.5,ymin=1e-11,ymax=1e-1,
  grid=both,minor tick num=1,
  legend style={font=\scriptsize,at={(0.5,-0.22)},anchor=north,
    legend columns=2,draw=none,fill=none},
  tick label style={font=\scriptsize},label style={font=\small},
  axis line style={layergrey},tick style={layergrey}]
  \addplot[layerindigo,thick,mark=*,mark size=1.7pt] coordinates {
    (0.4,0.01958) (0.2,0.009423) (0.1,0.004630)
    (0.05,0.002296) (0.025,0.001143)};
  \addplot[accent,thick,mark=square*,mark size=1.6pt] coordinates {
    (0.4,0.00001245) (0.2,0.0000007237) (0.1,0.00000004359)
    (0.05,0.000000002674) (0.025,0.0000000001656)};
  \legend{Euler (slope $1$), Runge--Kutta 4 (slope $4$)}
\end{axis}
\end{tikzpicture}
\caption{The error at $t=2$ for the logistic model as a function of the step $h$, on a
logarithmic scale on both axes. The Euler line falls with slope one: the error is
proportional to $h$. The Runge--Kutta line falls with slope four:
the error is proportional to $h^{4}$. At $h=0.4$ the two methods are about three
orders apart, and that distance widens when $h$ is made smaller. The numbers in the figure are
the same numbers as in the table in the text.}
\label{fig:galat-numerik}
\end{figure}

Now the property that deceives many people: stiffness. A system
is called \emph{stiff} if inside it there is a process that moves
far faster than the process we are interested in. Its practical signature: explicit methods
demand a very small step in order to remain \emph{stable}, even when the accuracy we ask for
demands only a large step.

The example is small and honest: $y'=-50(y-\cos t)-\sin t$, $y(0)=1$, whose exact
solution is $y=\cos t$. Here $y$ follows $\cos t$ slowly, but every
deviation is dragged back at the rate $50$ per second. Euler is stable only if
$|1-50h|<1$, that is $h<0.04$.

Results at $t=3$, with exact value $\cos3=-0.989992$: Euler with $h=0.02$
gives an error of $1.98\times10^{-4}$, with $h=0.039$ it gives
$7.97\times10^{-4}$, and with $h=0.05$ it gives $-1.839\times10^{7}$. One
step up from $0.039$ to $0.05$ changes the answer from almost right to
meaningless. RK4 holds out longer, because its stability region is wider: at
$h=0.05$ its error is $1.11\times10^{-3}$, but at $h=0.1$ it gives
$-1.32\times10^{31}$; its practical limit is $2.785/50=0.0557$.

The lesson to carry away: error and stability are two different things. The local error on
the first step is small and warns us of nothing; what explodes is
stability, and that is a matter of eigenvalues.

The antidote is implicit methods, the simplest being backward Euler:
\[
y_{n+1}=y_{n}+h\,f(t_{n+1},y_{n+1}),
\]
which demands solving an equation at every step, but is stable for large $h$
on stiff equations; for the example above, backward Euler with $h=0.5$
gives an answer close to $\cos t$. Reference: Hairer and Wanner, chapter IV.

\begin{tantangan}
Linearization gives a center, but a nonlinear system need not be a center. Construct a
two-dimensional system whose origin is \emph{not} a center even though its linearization is a
center, determine the character of that point (stable, asymptotically stable, or unstable),
and include an explicit formula, the eigenvalues of the Jacobian at the origin, and a brief
proof of your claim.

Hint 1: linearization sees only terms of degree one. Add a degree-three term that is
symmetric under rotation, then write the system in polar coordinates.

Hint 2: in polar coordinates, the rotation matrix becomes $\theta'=1$, and the radial
disturbance becomes $r'=cr^{3}$ for some constant $c$. The sign of $c$
decides everything, and $V=r^{2}$ is a Lyapunov function that works for
one of the two signs.

If you get stuck, look again at the example $r'=r^{3}$ in the Stability section.
\end{tantangan}
\section{Worked examples}
\tingkat{intermediate}

The four computations below are the backbone of this chapter; all of them can be checked
again with a calculator or with the program in the lab section, and all of them have gone
wrong in my own hands.

\textbf{Logistic model: one formula, two methods.}
Take $y'=y(1-y)$ with $y(0)=1/2$ and ask for the value at $t=2$. The formula from
the separation of variables section gives $y(2)=1/(1+e^{-2})=0.8807970780$. With
$h=0.1$ there are twenty steps: Euler gives $0.8854273272$, an error of
$4.630\times10^{-3}$ or $0.53$ percent, whereas Runge--Kutta with the same step
gives $0.8807970344$, an error of $4.359\times10^{-8}$. The cost is four times
as much, the accuracy about one hundred thousand times.

Note also the direction of the error. Euler is too high here, because
$y(1-y)$ is concave down on that interval, so a straight line always lies above the curve.
Checking the sign of the error is a cheap way to find out whether our program
works as it should.

\textbf{Variable coefficient: $y'+y/t=t^{2}$, $y(1)=0$.}
Here $p(t)=1/t$, so
\[
\mu(t)=\exp\Bigl(\int\frac{dt}{t}\Bigr)=t
\]
on $t>0$. Multiplying the equation by $t$ gives
\[
t\,y'+y=t^{3}
\quad\Longleftrightarrow\quad
\bigl(t\,y\bigr)'=t^{3},
\]
so $t\,y=t^{4}/4+C$ and
\[
y(t)=\frac{t^{3}}{4}+\frac{C}{t}.
\]
The condition $y(1)=0$ gives $C=-1/4$, so
\[
y(t)=\frac{t^{3}}{4}-\frac{1}{4t}=\frac{t^{4}-1}{4t},
\qquad y(2)=\frac{15}{8}=1.875 .
\]
Check: $y'=3t^{2}/4+1/(4t^{2})$ and $y/t=t^{2}/4-1/(4t^{2})$, their sum is
$t^{2}$, exactly the right-hand side. An adaptive integrator gives $1.875000000000$.

Without $\mu=t$, the left-hand side $y'+y/t$ does not look like the derivative of anything;
after multiplication by $t$, it becomes the derivative of $ty$.

\textbf{How large a swing can be before the small formula misses by one percent.}
The ratio of the exact period to the small period is
\[
\frac{T(\theta_{0})}{T_{0}}
=\frac{2}{\pi}K\Bigl(\sin^{2}\frac{\theta_{0}}{2}\Bigr),
\]
with $K$ the complete elliptic integral. Some of its values:

\begin{center}
{\small
\begin{tabular}{rrrr}
\hline
$\theta_{0}$ & $T/T_{0}$ & error & remark \\
\hline
$5^{\circ}$ & $1.000476$ & $0.048$ percent & pendulum clock \\
$10^{\circ}$ & $1.001907$ & $0.191$ percent & still safe \\
$20^{\circ}$ & $1.007669$ & $0.767$ percent & nearing the limit \\
$22.81^{\circ}$ & $1.010000$ & $1.000$ percent & the one-percent limit \\
$30^{\circ}$ & $1.017409$ & $1.741$ percent & correction needed \\
$60^{\circ}$ & $1.073182$ & $7.318$ percent & clock off by an hour \\
\hline
\end{tabular}}
\end{center}

So the small formula is fit to be used up to about twenty degrees, and
$22.81^{\circ}$ is the one-percent error threshold; the error grows like the square of the
amplitude, not like the amplitude.
Doubling the amplitude from $10^{\circ}$ to $20^{\circ}$ raises the error from
$0.19$ to $0.77$ percent, that is four times, in keeping with
$1+(\theta_{0}^{2}/16)$ as the first two terms of that estimate.

\kertas{
Test yourself in three steps. (1) Work the logistic example again with
$h=0.2$; you should get an Euler error of $9.423\times10^{-3}$ and an RK4 error of
$7.237\times10^{-7}$. (2) Derive the elliptic period formula yourself from energy
conservation, then check that $K$ at the parameter $0.25$ has the value
$1.6857503548$. (3) Find the threshold $h$ for Euler on $y'=-100y$; the answer is
$h<0.02$, and at $h=0.02$ exactly, the method oscillates between $+1$ and $-1$
without either decaying or exploding.}

\section{Computational lab}
\tingkat{advanced}

The following code computes all the numbers above: the Euler and RK4 errors on the logistic
model, the pendulum period by two different roads, and the stiffness threshold. It
uses only NumPy and SciPy.

\begin{lab}
\begin{verbatim}
import numpy as np
from scipy import integrate, special

def euler(f, y0, T, h):
    y, t, n = y0, 0.0, int(round(T / h))
    for _ in range(n):
        y += h * f(t, y); t += h
    return y

def rk4(f, y0, T, h):
    y, t, n = y0, 0.0, int(round(T / h))
    for _ in range(n):
        k1 = f(t, y); k2 = f(t + h/2, y + h*k1/2)
        k3 = f(t + h/2, y + h*k2/2); k4 = f(t + h, y + h*k3)
        y += h * (k1 + 2*k2 + 2*k3 + k4) / 6; t += h
    return y

logi = lambda t, y: y * (1 - y)
ex = lambda t: 1 / (1 + np.exp(-t))
print("logistic y(2) exact %.10f" % ex(2.0))
for h in (0.4, 0.2, 0.1, 0.05, 0.025):
    ee = abs(euler(logi, 0.5, 2, h) - ex(2))
    er = abs(rk4(logi, 0.5, 2, h) - ex(2))
    print("h=%.3f euler %.3e rk4 %.3e" % (h, ee, er))

g, L = 9.80665, 1.0
T0 = 2 * np.pi * np.sqrt(L / g)
th0 = np.radians(60.0)
k = np.sin(th0 / 2)
Tq = 4 * np.sqrt(L / (2*g)) * integrate.quad(
        lambda th: 1/np.sqrt(np.cos(th) - np.cos(th0)), 0, th0, limit=200)[0]
Te = 4 * np.sqrt(L / g) * special.ellipk(k**2)
print("T0 %.6f T60quad %.6f T60elliptic %.6f" % (T0, Tq, Te))
print("ratio %.6f percent %.4f" % (Te / T0, 100 * (Te / T0 - 1)))

stiff = lambda t, y: -50 * (y - np.cos(t)) - np.sin(t)
for h in (0.02, 0.039, 0.05):
    print("stiff euler h=%.3f y(3)=%.3e" % (h, euler(stiff, 1.0, 3, h)))
for h in (0.05, 0.1):
    print("stiff rk4 h=%.3f y(3)=%.3e" % (h, rk4(stiff, 1.0, 3, h)))
print("exact cos(3) %.10f" % np.cos(3))
\end{verbatim}
\end{lab}

The actual output, verbatim:

\begin{verbatim}
logistic y(2) exact 0.8807970780
h=0.400 euler 1.958e-02 rk4 1.245e-05
h=0.200 euler 9.423e-03 rk4 7.237e-07
h=0.100 euler 4.630e-03 rk4 4.359e-08
h=0.050 euler 2.296e-03 rk4 2.674e-09
h=0.025 euler 1.143e-03 rk4 1.656e-10
T0 2.006409 T60quad 2.153242 T60elliptic 2.153242
ratio 1.073182 percent 7.3182
stiff euler h=0.020 y(3)=-9.902e-01
stiff euler h=0.039 y(3)=-9.908e-01
stiff euler h=0.050 y(3)=-1.839e+07
stiff rk4 h=0.050 y(3)=-9.883e-01
stiff rk4 h=0.100 y(3)=-1.317e+31
exact cos(3) -0.9899924966
\end{verbatim}

Three things deserve reading. First, the Euler error ratios $2.08$, $2.04$, $2.02$,
$2.01$ and RK4 heading to $16=2^{4}$: first and fourth order are visible as numbers. Second,
the two roads to the pendulum period meet to six digits, even though quadrature and the
elliptic integral do not use the same formula. Third, the stiff lines: $h=0.039$
is still reasonable, $h=0.05$ is meaningless, and not one of those lines warns
us.

The program does not check the hypotheses of the theorems: whether $f$ satisfies the Lipschitz
condition, whether the equilibrium point is hyperbolic. That remains our job.

\begin{untukdosen}
The teaching order I use: start from the pendulum clock and the $7.32$
percent difference, then the definition of a solution along with the blow-up example
$1/(1-t)$, then separation of variables and the integrating factor. The integrating factor is
best proved in full at the board, because its proof uses only the product rule and the
fundamental theorem of calculus, two tools they have had since Chapter 13. Picard--Lindelof is
enough as a sketch with one iteration carried out; what must remain is the meaning of the
Lipschitz condition, not the fixed point technique. Phase plane before eigenvalues:
draw first, compute later.

Two misconceptions almost always appear. First, students take
``solution'' to mean a closed formula, then conclude that an equation without a formula has no
solution; the example $y'=\sqrt{|y|}$ helps, because there two solutions exist and
no single formula does. Second, they treat a small error as proof of correctness, and are
surprised when Euler with $h=0.05$ gives $-1.8\times10^{7}$.
Show them that output line; do not merely warn them about it.

One blackboard activity: draw the pendulum phase plane empty, then ask
the students to mark where the velocity is zero, where the acceleration is zero, and where the
energy is highest. The class will argue about the region between the two saddles, and that
argument is exactly the content of the theorem. A sign of readiness to go further: students
can guess the type of an equilibrium point from the trace and the determinant without
computing the roots.

\end{untukdosen}

\section{Summary}
\begin{ringkasan}
\begin{itemize}
\item A differential equation sets the rate of change equal to the state; its solution is a
function on an interval, and that interval is part of the answer.
\item Two families are solvable by hand: separation of variables and the first-order linear
equation with integrating factor $\mu=\exp(\int p\,dt)$.
\item Picard--Lindel\"of guarantees a unique solution on a short interval if $f$ is
continuous and satisfies the Lipschitz condition with respect to $y$; without that condition,
one initial state can give birth to two histories.
\item The phase plane replaces time with a picture: equilibrium points, closed trajectories,
separatrices, and open trajectories are read without any formula for the solution.
\item Local stability is read from the eigenvalues of the Jacobian, except when an eigenvalue
lies on the imaginary axis; there linearization does not decide, and a Lyapunov function is
used.
\item Euler is first order globally and RK4 is fourth order; both are explicit and therefore
fail on stiff systems, where stability, not accuracy, sets the size of the step.
\end{itemize}
\end{ringkasan}

\section{Exercises}
\begin{latihan}
\begin{enumerate}
\item Solve $y'=2ty$, $y(0)=1$, then compute $y(1)$ and compare it with
  $e$. Also compute Euler with $h=0.1$ and state its error. (My result:
  $2.33463336$, error $3.84\times10^{-1}$.)
\item Prove that $y'+p(t)y=q(t)$ with $p,q$ continuous has at most
  one solution, without an explicit formula: the difference of two solutions satisfies
  $w'=-p(t)w$ and $w(t_{0})=0$.
\item True, false, or not enough information: ``because linearization at the origin gives
  eigenvalues $\pm i$, the origin is a stable center''. Give one supporting example
  or one counterexample.
\item Prove that $E=v^{2}/2+\omega_{0}^{2}\theta^{2}/2$ is constant for the harmonic
  oscillator, then explain why its trajectories are ellipses and its period is
  $2\pi/\omega_{0}$.
\item An RC circuit with $R=2.2\,\mathrm{k}\Omega$, $C=10\,\mu\mathrm{F}$,
  $V_{s}=12$ volts starts from zero. Compute $\tau$, the voltage at $t=\tau$, and
  the time to reach $6$ volts. ($22$ milliseconds; $7.585447$ volts; $15.249$ milliseconds.)
\item Show that $v=\pm2\cos(\theta/2)$ satisfies $E=2$, then explain
  why that curve joins $(-\pi,0)$ to $(\pi,0)$ and is called a
  separatrix.
\item A pendulum with $L=0.25$ meters is released from $30^{\circ}$. Compute the small
  period and the exact period, then state the percentage difference. ($1.003205$ seconds;
  $1.020669$ seconds; $1.74$ percent.)
\item For $y'=-100y$ with $y(0)=1$, determine the threshold $h$ that keeps explicit Euler
  stable, then explain what happens exactly at that threshold.
  Check it with the program.
\item Prove that the scalar autonomous equation $y'=f(y)$ with $f$ continuous and
  satisfying the Lipschitz condition cannot have a nonconstant periodic solution.
\end{enumerate}

\textbf{Hints for the three hardest items.} Item 3: linearization sees only terms
of degree one; add a symmetric cubic term as in the example $r'=r^{3}$
and compare it with $r'=-r^{3}$. Item 6: use
$\cos\theta=2\cos^{2}(\theta/2)-1$ and check the value of $v$ at $\theta=\pm\pi$.
Item 9: suppose $y$ is periodic and nonconstant; at a maximum $y'=0$, so
$f(y_{\max})=0$, and uniqueness forces that trajectory to be constant.
\end{latihan}

\begin{untukpeneliti}
What is genuinely proved in this chapter, with written proofs: the integrating factor
formula together with its uniqueness, the complete Picard--Lindelof sketch with the error
bounds of the iteration, uniqueness of solutions of the linear equation through $w'=-pw$, the
Lyapunov test for both of its statements, and all the energy identities for the oscillator
and the pendulum. What I cite without proof: the Hartman--Grobman theorem (Hartman,
\emph{Ordinary Differential Equations}, chapter 9; Perko, \emph{Differential
Equations and Dynamical Systems}, section 2.8), Peano's theorem and
Arzel\`a--Ascoli, and the LaSalle invariance principle.

What is overpopularized is the phrase ``the eigenvalues decide everything''. What is true:
eigenvalues determine local behavior only when none of them has real part zero; on the
imaginary axis, degree-three terms can turn a center into a spiral, as in the example
$r'=r^{3}$ here. For high-dimensional systems, local ``asymptotic stability'' does not even
mean safe from large disturbances, because the basin of attraction can be very small.

An open question: how far can qualitative properties be certified numerically? A small error
on a single trajectory says nothing about stability; Euler with $h=0.039$ is right to four
digits, whereas $h=0.05$ gives $10^{7}$. For a system known only through samples of its
field, there is no general way to make sure that no stiff region hides between two sample
points.

Limits of validity: all functions here are real and smooth enough; all numbers are from
Python 3.14.4, NumPy 2.5.2, and SciPy 1.18.0, with quadrature at tolerance
$10^{-13}$ unless stated otherwise; $g=9.80665\ \mathrm{m/s^{2}}$ and friction is
neglected in the period computations.
\end{untukpeneliti}

\peta{Chapter 16 in Boyce and DiPrima, \emph{Elementary Differential Equations
and Boundary Value Problems}, chapters 1 (introduction), 2 (first-order equations), 3
(second-order linear equations), 4 (first-order systems and the phase plane), and 9
(nonlinear equations and stability); in Strogatz, \emph{Nonlinear Dynamics
and Chaos}, chapters 2 (flows on the line), 5 (two-dimensional systems), 6 (phase planes),
7 (oscillators), and 8 (bifurcations); in Arnold, \emph{Ordinary Differential
Equations}, chapters 1--4 for existence, uniqueness, and the classification of equilibrium
points; in Coddington and Levinson, chapters 1--2 for Picard--Lindelof and
Peano; in Hartman, chapter 9, and Perko, section 2.8, for Hartman--Grobman;
in Hairer and Wanner, \emph{Solving Ordinary Differential Equations II}, chapter
IV, for stiffness and implicit methods. The basis for the computations is the chapter
\emph{Derivatives: One and Several Variables} (partial derivatives and the Jacobian), the
chapter \emph{Series of Functions and Exchange of Limits} (uniform limits, which make Picard
iteration converge), and the chapter \emph{Multiple Integrals} (vector fields) in the first
edition of this book; the eigenvalues of Chapter 9 are used directly without being repeated,
and the floating-point error analysis of Chapter 22 is the continuation of the numerical part
of this chapter.}
